\documentclass[11pt, a4paper]{article}
\usepackage[utf8]{inputenc}
\usepackage{geometry}
\geometry{margin=1in}
\usepackage{hyperref}
\usepackage{listings}
\usepackage{xcolor}
\definecolor{codegray}{rgb}{0.5,0.5,0.5}
\definecolor{backcolour}{rgb}{0.95,0.95,0.92}
\lstdefinestyle{mystyle}{
    backgroundcolor=\color{backcolour},
    commentstyle=\color{codegray},
    keywordstyle=\color{blue},
    basicstyle=\ttfamily\footnotesize,
    breakatwhitespace=false,
    breaklines=true,
    captionpos=b,
    keepspaces=true,
    showspaces=false,
    showstringspaces=false,
    showtabs=false,
    tabsize=2
}
\lstset{style=mystyle}
\begin{document}
\begin{center}
    \large\bfseries Electron Momentum Transfer: User Manual
\end{center}
\section{Compilation}
The main program is compiled using CMake. It will automatically detect the host compiler and link the necessary libraries.
To compile the C binaries from source, execute the following commands in your terminal:
\begin{lstlisting}[language=bash]
git clone --depth 1 https://github.com/BanuDarius/electron-momentum-transfer
cd electron-momentum-transfer/
cmake -B build
cmake --build build
\end{lstlisting}
By default, CMake applies the \texttt{-march=native} and \texttt{-O3} optimization flags for your specific CPU architecture. If you are compiling on a legacy system or encounter illegal instruction errors, you can use CMake to build for a generic x86 architecture:
\begin{lstlisting}[language=bash]
cmake -B build -DGENERIC=ON
cmake --build build
\end{lstlisting}
To delete the previously generated output data, including images and video renders, use the clean target:
\begin{lstlisting}[language=bash]
cmake --build build --target clean-output
\end{lstlisting}
\section{Quick functions}
To easily check the functions of the program, you can uncomment one of the lines corresponding to these functions.
To run a low resolution parameter sweep:
\begin{lstlisting}[language=Python]
quick_example.run_quick_example(thread_num)
\end{lstlisting}
\begin{itemize}
    \item \textbf{\texttt{thread\_num}}: The number of threads to allocate for the simulation.
\end{itemize}
To reproduce one of the parameters sweeps presented in the \texttt{examples} directory:
\begin{lstlisting}[language=Python]
examples.run_example(num_example, thread_num)
\end{lstlisting}
\begin{itemize}
    \item \textbf{\texttt{num\_example}}: The  example to reproduce (1 to 4).
    \item \textbf{\texttt{thread\_num}}: The number of threads.
\end{itemize}
To run a complete test between the numerical integrator and an analytical solution for one plane wave:
\begin{lstlisting}[language=Python]
analytic_test.run_complete_test(thread_num)
\end{lstlisting}
\begin{itemize}
    \item \textbf{\texttt{thread\_num}}: The number of threads.
\end{itemize}
To replicate the results obtained by D. Bauer et al. (1995) Physical Review Letters paper:
\begin{lstlisting}[language=Python]
examples.replicate_prl_results(thread_num)
\end{lstlisting}
\begin{itemize}
    \item \textbf{\texttt{thread\_num}}: The number of threads.
\end{itemize}
To run a strong scaling performance test, from 1 to N threads:
\begin{lstlisting}[language=Python]
performance_test.run_example_performance_test(thread_num)
\end{lstlisting}
\begin{itemize}
    \item \textbf{\texttt{thread\_num\_final}}: The maximum number of threads for the test.
\end{itemize}
\section{Usage: \texttt{auto\_compute.py}}
The intended usage of this software suite is through the master Python script, \texttt{auto\_compute.py}. This script serves as the primary user interface and handles the complete simulation.\\
To start full parameter sweep and auto-generate the analysis figures, simply execute:
\begin{lstlisting}[language=bash]
python3 auto_compute.py
\end{lstlisting}
Configuration of the physical parameters (e.g., $a_0$, $\omega$, thread count) should be performed directly within the variables defined at the top of the \texttt{auto\_compute.py} script. The program uses atomic units.
\section{Laser Initialization and Parameters}
In the \texttt{auto\_compute.py} master script, the laser assembly is constructed by appending individual \texttt{LaserParameters} objects (defined in \texttt{sim\_init.py}) to the \texttt{lasers} list.
\begin{lstlisting}[language=Python]
lasers.append(sim_init.LaserParameters(a0, sigma, omega, etaf, zetax, zetay, alpha, phi, theta, psi, pond_integrate_steps, use_gaussian, w0, zeta_x_gauss, zeta_y_gauss))
\end{lstlisting}
The definitions of these arguments are:
\begin{itemize}
    \item \texttt{a0} ($a_0$): The dimensionless potential amplitude $a_0 = \frac{|e|E_0}{mc\omega}$.
    \item \texttt{sigma} ($\sigma$): The width parameter of the Gaussian temporal envelope.
    \item \texttt{omega} ($\omega$): The angular frequency of the laser field in atomic units.
    \item \texttt{etaf} ($\eta_f$): The flat-top envelope parameter. It defines the duration of the constant-amplitude plateau of the pulse in terms of the phase $\eta = \omega t - \mathbf{k} \cdot \mathbf{r}$. For a standard Gaussian pulse, set $\eta_f = 0$.
    \item \texttt{zetax}, \texttt{zetay} ($\zeta_x, \zeta_y$): The polarization parameters ($\mathbf{\epsilon}_1$ and $\mathbf{\epsilon}_2$). For instance, $\zeta_x = 0, \zeta_y = 1$ is for linear polarization, while $\zeta_x = \frac{1}{\sqrt{2}}, \zeta_y = \frac{1}{\sqrt{2}}$ is for circular polarization.
    \item \texttt{alpha} ($\alpha$): The polarization rotation angle. It rotates the basis vectors $\mathbf{\epsilon}_1$ and $\mathbf{\epsilon}_2$ around the propagation axis $\mathbf{n}$ by the angle $\alpha$.
    \item \texttt{phi}, \texttt{theta} ($\phi, \theta$): The azimuthal and polar spherical angles defining the laser's propagation direction vector $\mathbf{n}$.
    \item \texttt{psi} ($\psi$): The carrier-envelope phase (CEP), which shifts the phase of the oscillating carrier wave relative to the peak of the field envelope.
    \item \texttt{pond\_integrate\_steps}: A parameter used by the ponderomotive solver. It defines the number of steps used to evaluate the integrals (e.g., $\int A^2 d\phi$) with Simpson's 1/3 rule.
    \item \texttt{use\_gaussian}: A boolean which uses the Gaussian mode if true, and plane waves if false.
    \item \texttt{zeta\_x\_gauss, zeta\_y\_gauss}: Complex number polarization used in the Gaussian mode.
\end{itemize}\par
The line which runs the simulation itself is:
\begin{lstlisting}[language=Python]
programs.run_simulation(method, sim_parameters, lasers)
\end{lstlisting}
\begin{itemize}
    \item \textbf{\texttt{method}}: Defines the numerical integrator and physical approximation to be used. The Python string is internally mapped to a discrete \texttt{mode} integer passed to the C binary:
    \begin{itemize}
        \item \texttt{"electromagnetic"} $\rightarrow$ \texttt{mode = 0}: Uses the Higuera-Cary particle pusher for the full Lorentz force equations.
        \item \texttt{"ponderomotive"} $\rightarrow$ \texttt{mode = 1}: Uses the 4th-order Runge-Kutta (RK4) integrator to solve the averaged ponderomotive approximation.
        \item \texttt{"electromagnetic-rk4"} $\rightarrow$ \texttt{mode = 2}: Uses the RK4 integrator for the Lorentz force (used for legacy comparisons or specific debugging).
    \end{itemize}
    \item \textbf{\texttt{sim\_parameters}}: An instance of the \texttt{SimParameters} class.
    \item \textbf{\texttt{lasers}}: The list of \texttt{LaserParameters} objects.
\end{itemize}
\section{Error Analysis and Visualization}
The software is able to automatically benchmark the ponderomotive approximation against the full electromagnetic Lorentz force solver. 
This is executed with:
\begin{lstlisting}[language=Python]
programs.calculate_errors(sim_parameters)
\end{lstlisting}
\begin{itemize}
    \item \textbf{\texttt{sim\_parameters}}: The current configuration state.
\end{itemize}
The average errors are plotted with:
\begin{lstlisting}[language=Python]
plotting.plot_average_errors_all(a0_array)
\end{lstlisting}
\begin{itemize}
    \item \textbf{\texttt{a0\_array}}: The linear space of $a_0$ values swept during the simulation.
\end{itemize}
The heatmaps for visualizing the final momentum as a function of $a_0$:
\begin{lstlisting}[language=Python]
plotting.plot_2d_heatmap_all(method, sim_parameters, a0_array, axis_pos)
\end{lstlisting}
\begin{itemize}
    \item \textbf{\texttt{method}}:\texttt{"electromagnetic"} or \texttt{"ponderomotive"}.
    \item \textbf{\texttt{sim\_parameters}}: The current simulation state.
    \item \textbf{\texttt{a0\_array}}: The linear space of $a_0$ values swept during the simulation.
    \item \textbf{\texttt{axis\_pos}}: The coordinate axis index corresponding to the initial distribution of the particles (the X-axis of the heatmap).
\end{itemize}
To plot the errors between the two methods:
\begin{lstlisting}[language=Python]
plotting.plot_2d_errors_heatmap_all(sim_parameters, a0_array, axis_pos)
\end{lstlisting}
\begin{itemize}
    \item \textbf{\texttt{sim\_parameters}}: The current simulation state.
    \item \textbf{\texttt{a0\_array}}: The linear space of $a_0$ values swept during the simulation.
    \item \textbf{\texttt{axis\_pos}}: The coordinate axis index corresponding to the initial position of the electrons.
\end{itemize}
This function automatically performs a numerical convergence test by re-running the simulation with a strictly higher temporal resolution and comparing the deviation in the final momentum.
\begin{lstlisting}[language=Python]
programs.check_convergence(method, sim_parameters, lasers, multiplier)
\end{lstlisting}
\begin{itemize}
    \item \textbf{\texttt{method}}: The physics solver to validate.
    \item \textbf{\texttt{sim\_parameters}}: The current simulation configuration.
    \item \textbf{\texttt{lasers}}: The list of laser parameters.
    \item \textbf{\texttt{multiplier}}: The factor by which to increase the number of integration steps.
\end{itemize}
This function generates a plot of the average convergence error as a function of $a_0$:
\begin{lstlisting}[language=Python]
plotting.plot_convergence_all(method, a0_array)
\end{lstlisting}
\begin{itemize}
    \item \textbf{\texttt{method}}: Specifies the physics solver utilized.
    \item \textbf{\texttt{a0\_array}}: The linear space of $a_0$ values swept during the simulation (mapped to the X-axis).
\end{itemize}
To generate the complete convergence heatmap:
\begin{lstlisting}[language=Python]
plotting.plot_2d_convergence_heatmap_all(method, sim_parameters, a0_array)
\end{lstlisting}
\begin{itemize}
    \item \textbf{\texttt{method}}: Specifies the physics solver being validated.
    \item \textbf{\texttt{sim\_parameters}}: The current simulation configuration.
    \item \textbf{\texttt{a0\_array}}: The linear space of $a_0$ values swept during the simulation (the Y-axis).
\end{itemize}
\end{document}